\paragraph{Conjugate-coordinate class times a fuzzy factor: torsion of the level-one amplitudes and the induced Einstein coefficient (conditional; bounded no-go and conditional route).}
This paragraph supplements the preceding one, in which the two-derivative density of the mixed term on the six level-one amplitudes, its cone values and the torsion those amplitudes source were not computed. Take for \(M\) the conjugate-coordinate class (Poisson limit, leaf \(\mathbb{CP}^{1,2}\), \(X^a=T^a+A^a\), \(T^a=M^{a4}\), base functions \(y^\mu=M^{\mu5}\), \(m=M^{45}\)), the split \(d=4\) and \(K\) of the preceding paragraph, the level-one amplitudes \(A^a=c_i^a(y)\,n^i\) with \(M^{\mu\rho}=\sum_in^iS_i^{\mu\rho}\) on the fibre, \(n\) a unit vector of the fibre two-sphere, a plane wave \(e^{\mathrm iq\cdot y}\) with frozen coefficients at a base point, \(|q|\ell\gg1\), \(s_i^\mu=S_i^{\mu\nu}q_\nu\), \(G_{ij}=s_i\cdot s_j\), and the leading symbol in which the fibre momenta and the fibre components of the Hamiltonian fields are dropped and the trace over the fibre harmonics of a multiplication operator is its fibre average. This is the hybrid evaluation of Section~7 of~\cite{Steinacker2023}, for a finite fibre and a scale of \(K\) above the fibre fluxes.

\emph{Frame and torsion.} With \(E^{a\mu}=\{X^a,y^\mu\}\), \(F^{ab}=\{X^a,X^b\}\) and \(T^{ab\mu}=\{F^{ab},y^\mu\}\), the background has \(E^{a\mu}=\pm m\,\eta^{a\mu}\) and \(T^{ab\mu}=\eta^{a\mu}y^b-\eta^{b\mu}y^a\), independent of the fibre point (exact in the Lie--Poisson algebra). At the leading symbol, for an amplitude \(\alpha^a\) with any dependence on the fibre point, \(\delta E^{a\mu}=\alpha^a v^\mu\) up to a constant phase, with \(v=\sum_jn^js_j\), so the frame perturbation at a fibre point has rank one, and \(\delta T^{ab\mu}=-m\,v^\mu(q^a\alpha^b-q^b\alpha^a)\); for the level-one amplitudes \(\alpha=\sum_in^ic_i\) (verified by a Lie--Poisson computation at points of an explicit realisation of the orbit, relative residual \(3\cdot10^{-14}\)). Write \(c_i=D_{ij}s_j+\gamma_iq\). The \(\gamma_i\) are gauge directions, the fibre-averaged metric depends only on the symmetric part of \(D\), which gives the six polarisations, and the antisymmetric part gives the three inert amplitudes. The torsion of the fibre-averaged frame has no axial part on the six polarisations and a vector part proportional to \(q^a\operatorname{tr}(DG)\), that is to the trace of the averaged metric perturbation; the inert amplitudes carry axial torsion. At a fibre point the six polarisations do source axial torsion, proportional to \([D,nn^{T}]\) and so in the second fibre harmonic; it averages to zero and vanishes identically only for \(D\propto\mathbb I\) (exact rational arithmetic in seven sampled cases: four base points, five wave vectors including a time-like one, both orientations of the fibre basis).

\emph{Two-derivative forms.} In the normalisation of the preceding paragraph, for a real wave averaged over its phase, the index density at a fibre point is \(X=m^{-4}(v\cdot v)\,(q^2\alpha\cdot\alpha-(q\cdot\alpha)^2)\), and the scalar density is \(\tfrac16\sqrt{|G|}R[G]-\sqrt{|G|}\,G^{\mu\nu}\partial_\mu\phi\,\partial_\nu\phi\), with \(G\) the inverse of \(\eta_{ab}E^aE^b\) and \(\phi=-\tfrac14\ln|\det G|\). With these densities the two modes with \(\delta\rho=0\) of four coordinates have the index form \(-4\) and the scalar form \(\tfrac16\) times the Einstein--Hilbert form, as in the preceding paragraph. For the fibre-averaged frame, exactly at every sampled point, the index form is \(-2\) times the Einstein--Hilbert form \(\mathrm{EH}(\bar h)\) of the averaged metric plus a multiple of \(q^2(\operatorname{tr}\bar h)^2\), and the scalar form is \(\tfrac16\,\mathrm{EH}(\bar h)\) plus such a multiple. The inert amplitudes do not couple to \(\bar h\) and enter only through the axial term. Hence, in the units \(c\) of the preceding paragraph, \(c_T=\tfrac16W_R-2W_X\) on the five traceless polarisations and \(-\tfrac5{24}W_R+W_X\) on the trace polarisation, their complement in the six. With the cone values of the preceding paragraph and the cone in the base at every fibre point, the fixed-state replica has \(c_A=\tfrac16W_R-2W_X\), equal to \(c_T\) on the five traceless polarisations for the finite parts, without the dressed density of the paragraph on local frame rotations. The periodic and smooth gaps there are \(2W_X\) and \(-\tfrac14W_R+2W_X\). On the trace polarisation the fixed-state, periodic and smooth gaps are \(\tfrac38W_R-3W_X\), \(\tfrac38W_R-W_X\) and \(\tfrac18W_R-W_X\).

\emph{Fibre trace.} The fibre averages of the pointwise densities differ from these by the second fibre harmonic. On the six polarisations neither form lies in the span of \(q^2\bar h_{\mu\nu}\bar h^{\mu\nu}\) and \(q^2(\operatorname{tr}\bar h)^2\) (relative residual of the best fit between \(0.03\) and \(0.36\)), both couple the inert amplitudes to the six, and their generalised eigenvalues relative to \(\mathrm{EH}(\bar h)\) on the five traceless polarisations are distinct and of both signs (for the index form \(-4.36\), \(-0.73\), each twice, and \(0.47\) at one point). No ratio \(W_R:W_X\) makes their sum proportional to \(\mathrm{EH}(\bar h)\) or covariant (exact linear algebra at the seven sampled points; a Monte Carlo fibre average on the explicit realisation agrees). For this trace no single \(G_T\) covers the six polarisations or the five traceless ones, and \(-2\) is not a generalised eigenvalue of the index form at any sampled point, so the mixed-term part of the fixed-state gap vanishes on no traceless polarisation. The trace with equal weight on all fibre harmonics applies when the scale of \(K\) exceeds the fibre gaps, which is the regime of the cited construction. The averaged-frame forms are the trace over the fibre-constant loop modes with projected coefficients. The full trace approaches them when the fibre gaps exceed the scale of \(K\), since intermediate states in the second harmonic are then suppressed in the proper-time integral; this is an argument and is not computed. In that regime the first term of the factorisation has a part with two base derivatives from the fibre flux, the mixed \(S^2\)--\(M\) term of Section~7.1.2 of~\cite{Steinacker2023}, which that work states is not Lorentz invariant a priori and finds subdominant for a scale of \(K\) above the fibre scale, a condition that fails in that regime; this part is not computed here. In the cited construction the mixed term is written with the torsion of a frame on \(M^{3,1}\). Replacing the frame of the level-one amplitudes by its fibre average gives the averaged-frame forms above, and for a fibre-independent torsion, such as that of the background, the two evaluations agree.

These statements are for the Poisson limit, one loop, the content and product backgrounds of the preceding paragraph, the orientation-even mixed term, the level-one amplitudes, the leading symbol at a frozen base point with \(y^i=0\) (other points follow by Lorentz covariance), a formal evaluation with the Lorentzian \(\eta\), and the replicas of Section~\ref{sec:status-B} with the cone in the base. They do not compute the first-term contributions of order \(\LUV^2\), the logarithmic and fibre-flux contributions, the subleading symbol, the interpolation between the two regimes, or the fibre spectrum. They do not construct \(K\) or fix the sign of \(W_R\), derive the replica or the regime from the matrix state, identify \(\bar h\) with the emergent metric of Chapter~\ref{ch:vierbein-emergence}, or bear on Assumption~\ref{ass:lorentzian}. Problems~B and~F remain open and PROGRAMME-CONDITIONAL.
